% GENERATED FILE. Edit canonical lessons, then run npm run generate:books.
% canonical-source-hash: fnv1a64:1073df674e617440
% canonical-lessons: FR-C430-punaise, FR-C430-cafard, FR-C430-moustiquaire, FR-C430-volet, FR-C430-store

\chapter{Bestioles et volets}
\label{ch:fr-bestioles-et-volets}
\hlchaptermodality{430}

\begin{chapteropening}
\textbf{By the end of this chapter:} \emph{I can complain about bugs and light in a room: the bedbug, the cockroach, the mosquito net, the shutter and the blind.}

Complete the last lesson of chapter 430: I can complain about bugs and light in a room: the bedbug, the cockroach, the mosquito net, the shutter and the blind.
\end{chapteropening}

\section[la punaise]{la punaise — the bedbug; the drawing pin}

\label{lesson:FR-C430-punaise}



\begin{quote}
Before the new one: say the French for the fitted carpet, then the French for the stain.
\end{quote}

\subsection*{You'll want to know: la punaise}
\textbf{la punaise} — \textquotedblleft{}the bedbug; the drawing pin\textquotedblright{}.

\begin{quote}
Elle a trouvé des punaises dans le matelas.
\end{quote}

\begin{quote}
\emph{She found bedbugs in the mattress.}
\end{quote}

\emph{Une punaise de lit} — a bedbug; in an office, \emph{une punaise} is a drawing pin.

\begin{cousinweb}
\emph{punaise} is built on the old adjective \emph{punais} 'stinking' (EtymDB), which goes back to Latin \emph{putēre} 'to stink'.
\end{cousinweb}

\begin{grammarlens}[title={What you've built}]
One more word for telling, explaining and describing at B1.
\end{grammarlens}

\subsection*{Your turn}
\begin{itemize}
\raggedright
  \item \emph{Say it:} \emph{la punaise}
  \item \emph{Say it:} \emph{la punaise}, once more
  \item \emph{Say it:} a sentence of your own with \emph{la punaise}
  \item \emph{Recall:} say the French for the fitted carpet, then the French for the stain, then say \emph{la punaise} again
\end{itemize}

\begin{tcolorbox}[breakable,colback=teal!4,colframe=teal!35!black,title={Before you move on}]
What does la punaise mean? (\textquotedblleft{}The bedbug; the drawing pin\textquotedblright{}.) Say it once more.
\end{tcolorbox}
\section[le cafard]{le cafard — the cockroach; the blues}

\label{lesson:FR-C430-cafard}



\begin{quote}
Before the new one: say the French for the stain, then the French for the bedbug.
\end{quote}

\subsection*{You'll want to know: le cafard}
\textbf{le cafard} — \textquotedblleft{}the cockroach; the blues\textquotedblright{}.

\begin{quote}
Il y avait un cafard dans la salle de bains.
\end{quote}

\begin{quote}
\emph{There was a cockroach in the bathroom.}
\end{quote}

\emph{Avoir le cafard} — to feel low.

\begin{cousinweb}
\emph{cafard} is usually traced to Arabic \emph{kāfir} 'unbeliever', first an insult for a two-faced person; the path to the insect is not settled. EtymDB lists the word without a source.
\end{cousinweb}

\begin{grammarlens}[title={What you've built}]
One more word for telling, explaining and describing at B1.
\end{grammarlens}

\subsection*{Your turn}
\begin{itemize}
\raggedright
  \item \emph{Say it:} \emph{le cafard}
  \item \emph{Say it:} \emph{le cafard}, once more
  \item \emph{Say it:} a sentence of your own with \emph{le cafard}
  \item \emph{Recall:} say the French for the stain, then the French for the bedbug, then say \emph{le cafard} again
\end{itemize}

\begin{tcolorbox}[breakable,colback=teal!4,colframe=teal!35!black,title={Before you move on}]
What does le cafard mean? (\textquotedblleft{}The cockroach; the blues\textquotedblright{}.) Say it once more.
\end{tcolorbox}
\section[la moustiquaire]{la moustiquaire — the mosquito net; the fly screen}

\label{lesson:FR-C430-moustiquaire}



\begin{quote}
Before the new one: say the French for the bedbug, then the French for the cockroach.
\end{quote}

\subsection*{You'll want to know: la moustiquaire}
\textbf{la moustiquaire} — \textquotedblleft{}the mosquito net; the fly screen\textquotedblright{}.

\begin{quote}
La fenêtre a une moustiquaire, alors nous dormons bien.
\end{quote}

\begin{quote}
\emph{The window has a fly screen, so we sleep well.}
\end{quote}

Built on \emph{le moustique}.

\begin{cousinweb}
\emph{moustiquaire} is built on \emph{moustique} 'mosquito', from Spanish \emph{mosquito} 'little fly', which goes back to Latin \emph{musca} 'fly' (Lewis \& Short).
\end{cousinweb}

\begin{grammarlens}[title={What you've built}]
One more word for telling, explaining and describing at B1.
\end{grammarlens}

\subsection*{Your turn}
\begin{itemize}
\raggedright
  \item \emph{Say it:} \emph{la moustiquaire}
  \item \emph{Say it:} \emph{la moustiquaire}, once more
  \item \emph{Say it:} a sentence of your own with \emph{la moustiquaire}
  \item \emph{Recall:} say the French for the bedbug, then the French for the cockroach, then say \emph{la moustiquaire} again
\end{itemize}

\begin{tcolorbox}[breakable,colback=teal!4,colframe=teal!35!black,title={Before you move on}]
What does la moustiquaire mean? (\textquotedblleft{}The mosquito net; the fly screen\textquotedblright{}.) Say it once more.
\end{tcolorbox}
\section[le volet]{le volet — the shutter}

\label{lesson:FR-C430-volet}



\begin{quote}
Before the new one: say the French for the cockroach, then the French for the mosquito net.
\end{quote}

\subsection*{You'll want to know: le volet}
\textbf{le volet} — \textquotedblleft{}the shutter\textquotedblright{}.

\begin{quote}
Le matin, elle ouvre les volets de la chambre.
\end{quote}

\begin{quote}
\emph{In the morning she opens the bedroom shutters.}
\end{quote}

\emph{Fermer les volets} — to close the shutters.

\begin{cousinweb}
\emph{volet} is built on \emph{voler} 'to fly' (EtymDB), which is inherited from Latin \emph{volāre}: a panel that swings open like a wing.
\end{cousinweb}

\begin{grammarlens}[title={What you've built}]
One more word for telling, explaining and describing at B1.
\end{grammarlens}

\subsection*{Your turn}
\begin{itemize}
\raggedright
  \item \emph{Say it:} \emph{le volet}
  \item \emph{Say it:} \emph{le volet}, once more
  \item \emph{Say it:} a sentence of your own with \emph{le volet}
  \item \emph{Recall:} say the French for the cockroach, then the French for the mosquito net, then say \emph{le volet} again
\end{itemize}

\begin{tcolorbox}[breakable,colback=teal!4,colframe=teal!35!black,title={Before you move on}]
What does le volet mean? (\textquotedblleft{}The shutter\textquotedblright{}.) Say it once more.
\end{tcolorbox}
\section[le store]{le store — the blind (on a window)}

\label{lesson:FR-C430-store}



\begin{quote}
Before the new one: say the French for the mosquito net, then the French for the shutter.
\end{quote}

\subsection*{You'll want to know: le store}
\textbf{le store} — \textquotedblleft{}the blind (on a window)\textquotedblright{}.

\begin{quote}
Le store est cassé, il ne descend plus.
\end{quote}

\begin{quote}
\emph{The blind is broken; it will not come down any more.}
\end{quote}

\emph{Baisser le store} — to pull down the blind.

\begin{cousinweb}
\emph{store} is usually traced to Italian \emph{stora}; behind it is Latin \emph{storea} 'mat of rushes' (EtymDB).
\end{cousinweb}

\begin{grammarlens}[title={What you've built}]
One more word for telling, explaining and describing at B1.
\end{grammarlens}

\subsection*{Your turn}
\begin{itemize}
\raggedright
  \item \emph{Say it:} \emph{le store}
  \item \emph{Say it:} \emph{le store}, once more
  \item \emph{Say it:} a sentence of your own with \emph{le store}
  \item \emph{Recall:} say the French for the mosquito net, then the French for the shutter, then say \emph{le store} again
\end{itemize}

\begin{tcolorbox}[breakable,colback=teal!4,colframe=teal!35!black,title={Before you move on}]
What does le store mean? (\textquotedblleft{}The blind (on a window)\textquotedblright{}.) Say it once more.
\end{tcolorbox}
